\section{Six integral properties of the level-one Wach module}

Let $S$ be a local situation in the top-degree theorem. Its standing
assumptions give $p\ge5$, an unramified field $K=K_0$, a good
level-one layer, and coefficients $O=\mathcal O_L$ splitting the
embeddings of $K$ and containing $\mu_p$. Use the very module $N(T)$
from L78W13 and F78D2. On its $i$th component write $A=O[[X]]$,
$Q=\Phi_p(1+X)$, and $P_i=A\varphi N_{i-1}$. The tame subgroup
$\Delta$ has cyclotomic characters $d$ ranging over $\mu_{p-1}$ in
$\mathbb Z_p^\times$. Indeed the cyclotomic tower is totally ramified
and hence disjoint from $K$ over $\mathbb Q_p$.

% BEGIN W80N1
\paragraph{W80N1. The integral invariant parameter and its primitive divisor.}
The source's parameter on each letter is
\[
 v=1-p+\sum_{d\in\mu_{p-1}}(1+X)^d.
\]
It is invariant and has constant term zero. For $1\le j<p-1$ the
coefficient of $X^j$ vanishes. The polynomial $\binom{d}{j}$ in $d$
has zero constant term and the sum of $d^a$ over $\mu_{p-1}$ vanishes
for $0<a<p-1$. For $j=p-1$ the leading coefficient gives
$(p-1)/(p-1)!$, a unit. Thus $v=X^{p-1}$ times a unit.
Set $Z=(p-1)^{-1}\sum_d d^{-1}((1+X)^d-1)$. It has linear term
$X$ and transforms under $d$ by multiplication by $d$. It is another
integral power series coordinate by successive coefficient inversion.
The invariants are $O[[Z^{p-1}]]$ by comparing coefficients; distinct
tame roots remain distinct in the residue field. In this ring $v$
has unit linear term in $Z^{p-1}$, hence $A^\Delta=O[[v]]$.
The same argument identifies the residual invariants with $k[[\bar v]]$.
At every primitive $p$th root $\zeta$, the sum of the $\zeta^d$ is
$-1$, hence $v(\zeta-1)=-p$. Integral division by each distinct linear
factor as in W79H4 gives $v+p=Qh$ with $h\in A$. At $X=0$ this gives
$p=ph(0)$, so $h$ is a unit. In particular $v+p$ is a uniformizer
of the primitive Taylor ring on every branch.
% END W80N1

% BEGIN W80N2
\paragraph{W80N2. An invariant basis, both height bounds and the signed rank-one equality.}
Start with the integral basis of L78W1. At $\zeta=1$ the expansion
basis generates $D_i[[Y]]$ by W79H3 and is a basis there (it spans
in the correct rank). Each $\delta\in\Delta$ acts in the expansion
by substitution $Y\mapsto dY$ with fixed crystalline constants:
$K=K_0$, so all the constants in the chosen crystalline basis are
fixed by $G_K$. Evaluation at zero therefore gives the identity
matrix on this basis. The action matrix lies in $\mathrm{GL}_r(A)$,
so its constant term is already the identity over $O$. Averaging
each basis vector over $\Delta$, dividing by $p-1$, changes its matrix
by a multiple of $X$. The resulting determinant has constant term one.
Doing this on each letter gives a common invariant basis of $N(T)$.
Its invariant coefficients are exactly the series in $v$ of W80N1.
W79H4 applied with the bounds on each embedding gives both native
height inclusions for all jump multisets. In rank one, if $a_i$ is
the jump, W79H2--W79H3 give $Y^{-a_i} D_i[[Y]]$ for the expansion
of $N_i$ and $D_i[[Y]]$ for that of $P_i$. Write $h^+=\max(h,0)$.
Then the primitive expansions of $Q^{(-a_i)^+}P_i$ and
$Q^{a_i^+}N_i$ coincide. Their basis changes have denominators only
at $Q$, by the original two Wach inclusions. W79H4's numerator
division at every root, applied in both directions, gives equality
of the two integral lattices. This includes negative and zero jumps.
% END W80N2

% BEGIN W80N3
\paragraph{W80N3. Isogenies reflect membership and integral direct sums split.}
In a Wach basis the inclusion $N\subset D(T)$ is coordinatewise
$A\subset\mathcal E_O$, the $p$-adic completion of $A[1/X]$.
Indeed the basis extends by the base change field of the Wach data.
If $p^m x\in A$ for $x\in\mathcal E_O$ then $x\in A$:
$A/p^m A\longrightarrow(A/p^m A)[1/X]$ is injective since
multiplication by $X$ just shifts coefficients. The completion has
this very reduction (the reductions are surjective, with $p$ regular),
and is $p$-torsion-free by its Laurent coefficient description.
Now a rational isomorphism $f:T'\to T$ preserving lattices has
$p^m f^{-1}$ integral for sufficiently large $m$, by finitely many
matrix entries. Crystallinity transports to $T'$. Apply the integral
map on Wach modules for $p^mf^{-1}$ to $D(f)x\in N(T)$, obtaining
$p^mx\in N(T')$ by the composition and scalar laws. The coordinate
test gives $x\in N(T')$, the asserted saturation. For the direct-sum
clause, the rational isomorphism $(f,g)$ is surjective on the given
integral lattices by hypothesis. Its inverse has integral coordinate
projections $u,w$. They commute with the group action. The target is
crystalline by the direct sum of invariant comparison bases.
Functoriality takes $u,w$ into their respective Wach lattices.
Applying the two inclusions back gives the original vector: the
identity $fu+gw=1$ holds after tensoring with the Witt comparison
ring, where the maps act on representation coordinates, and the
comparison inclusion is injective. This proves the native additivity.
% END W80N3

% BEGIN W80N4
\paragraph{W80N4. Divided cyclotomic action and its polynomial on all primitive jets.}
Let $\gamma$ have character $1+p$. On $N^\Delta$, $\gamma-1$ takes
values in $XN$ by the level-one action law. Its image is invariant
since the cyclotomic group is abelian. In the invariant basis of
W80N2 its coefficients thus belong to $O[[v]]$ and vanish at zero,
so are divisible by $v$. Put $\nu=(\gamma-1)/v$, an $O$-linear
endomorphism of $N^\Delta$. Let $c(j)=(1-(1+p)^j)/p$ in $O$ for
$j\in\mathbb Z$. This is also the only possible choice in the source's
normalization, by injectivity of $O\to L$.

At a primitive root choose a basis of $D_i$ adapted to the filtration.
By W79H2--W79H3 a line of jump $u$ with basis $b$ contributes
$Y^{-u}bL[[Y]]$ to the Taylor lattice. The operator on that line
sends $F(Y)b$ to $(F((1+p)Y)-F(Y))b/v(Y)$, since $\gamma$ fixes
$\zeta$ and fixes the crystalline basis. Here $v(0)=-p$ in the
primitive Taylor coordinate. On its quotient modulo $Y^s$ times
that lattice the matrix is triangular with diagonal entries
$c(k-u)$ for $0\le k<s$. These are distinct on this line, since
$1+p$ has infinite order. The product of their linear factors
annihilates the quotient by Cayley--Hamilton. Each of the lines is
preserved, including when two lines share a jump. The product over
the \emph{set} of all $k-u$ therefore annihilates the entire quotient.
This also works at $s=0$ where the quotient is zero. For invariant
$y$ the polynomial applied to $y$ lies integrally in $N$. Its
coordinates on the $i$th letter vanish to order at least $s$ at
every primitive root. W79H4's division test and W80N1 give divisibility
by $(v+p)^s$ in $N_i$ as required.
% END W80N4

% BEGIN W80N5
\paragraph{W80N5. Characteristic polynomial in every invariant basis and the full integral record.}
Take any invariant $A$-basis $b_j$ in the native universal clause.
The columns of $\nu b_j$ have coefficients $B_{j'j}(v)\in O[[v]]$
on each letter. Evaluation at $v=-p$ converges in $O$, producing a
matrix $B(-p)$ over $O$. The difference of each column from the
result at $-p$ is divisible by $v+p$ by integral power series
division. At a primitive root the change to W80N4's adapted basis
is invertible on the Taylor lattice. On its quotient of length one
$\nu$ is $L$-linear: substitution fixes constants and is identity
on coefficient series modulo $Y$. Its matrix in the evaluated
basis is $B(-p)$, and the adapted basis has eigenvalues $c(-u)$
with the multiplicities of the jumps. Thus its characteristic
polynomial is $\prod_{u\in U_i}(X-c(-u))$ in $O[X]$, by injectivity
into $L[X]$. W80N4 and this calculation prove both polynomial clauses
of $\mathrm{SenOp}$. With W80N2--W80N3 this supplies each of
$\mathrm{InvBasis}$, Heights, RankOne, Saturated, Additive and SenOp
for the original $\mathrm{QN}$ record on our chosen $S.W$.
% END W80N5
